\chapterimage{container.jpg}

\chapter{Això és teu?}\label{sec:packets}

La informació viatja per aquest graf de xarxes en forma de paquets, que contenen una petita quantitat 
d'informació i de metainformació. 

Els dos extrems de la comunicació han d'acordar el format dels paquets, de manera que aquests paquets es puguin produir 
i interpretar de manera eficient i automàtica.

Quan els paquets viatgen per una o més xarxes, es poden perdre o desordenar. Si entre dos ordinadors ha de circular 
un flux continu de dades, els paquets han de contenir mecanismes per detectar les pèrdues i restaurar l'ordre correcte.

\section{Tinc una entrega per a tu}\label{sec:packets:lan_wan}

No és viable establir connexions físiques entre cada parell d'ordinadors que es volen comunicar.
En lloc d'això, es comparteixen relativament poques connexions per transportar missatges entre molts parells de dispositius.

Si un sol dispositiu transmet dades de manera contínua durant molt de temps, aquella connexió queda bloquejada i pot esdevenir un coll d'ampolla que 
impedeix que altres dispositius es comuniquin. Per evitar aquest problema, les connexions limiten la quantitat màxima de dades
que es pot enviar sense interrupció. Com a resultat, els dispositius es veuen obligats a enviar dades en ràfegues discretes anomenades 
\conceptRef{packet}{paquets}. L'intercanvi d'aquest tipus de missatges a través d'una o més xarxes s'anomena
\conceptRef{packet switching}{commutació de paquets}.

\begin{remark}
% The concept of \concept{packet switching} (crucial to the creation of the Internet)
% was developed in the 1960's (during the cold war) by the US military
% to have a communications network that would work in case of a nuclear strike.
%  
% The problem with the one they had was that, if one station was damaged, the whole system stopped working.
% However, by using packet switching, any connected subgraph of the network would keep working.
% 
% The reasoning for creating a packet-switched network was that they could still retaliate their enemy in case of an attack, which would
% make them a less attractive target.

Leonard Kleinrock, un dels pioners de la interconnexió de xarxes, disputa l'autoria de la commutació de paquets,
que sovint s'atribueix a Paul Baran i Donald Davies.
\end{remark}

Els \conceptRef{packet}{paquets} han de ser petits perquè les connexions no quedin bloquejades durant massa temps. 
Alhora, volem que els paquets continguin tantes dades com sigui possible, perquè cada paquet 
conté dades de \conceptRef{overhead}{sobrecàrrega} (com ara les adreces que es tracten a la secció~\ref{sec:where}) que cal
enviar a més de les dades de payload (càrrega útil) de l'usuari. Les xarxes fixen la longitud màxima dels paquets
mitjançant un paràmetre anomenat \conceptRef{MTU}{unitat màxima de transmissió} (MTU), \eg, 1500~bytes a Ethernet.


La transmissió de paquets a través d'un \conceptRef{data link}{enllaç de dades} s'ha de fer amb cura,
sobretot quan aquest enllaç és un \concept{bus} compartit per diversos dispositius.
Cada paquet s'ha de poder distingir individualment de la resta, 
cosa que es pot fer amb almenys tres estratègies:
\begin{enumerate}
\item \textit{Longitud fixa}: tots els paquets tenen la mateixa longitud. L'emissor i el receptor han
d'haver acordat aquest valor per endavant.

\item \textit{Longitud explícita}: la longitud del paquet s'inclou dins del paquet, \eg, 
fent servir el primer byte del paquet. L'emissor i el receptor han d'acordar com s'inclou aquesta informació
i com s'ha d'interpretar.

\item \textit{Seqüències d'escapament}: certes \conceptRef{escape sequence}{seqüències d'escapament} binàries 
dins de les dades tenen un significat especial. 
Per exemple, es podria definir que el byte \otherBase{11110000} (\otherBase{0xf0}) 
significa ``final de paquet'': aquests bits han d'aparèixer després de cada paquet, i només aleshores.
% 
Llavors l'emissor podria començar a transmetre el contingut d'un paquet, seguit de \otherBase{0xf0}.
Les dues parts han d'acordar quines seqüències d'escapament fan servir, què signifiquen i com 
enviar dades iguals a aquestes seqüències (\eg, hauria de ser possible enviar \otherBase{0xf0f0f0}
sense provocar un ``final de paquet'' fins que s'hagin transmès tots els bytes).
\end{enumerate}

\begin{remark}
També és possible senyalitzar l'\textit{inici} d'un paquet. En aquest cas, es fa servir una seqüència prèvia anomenada \conceptRef{preamble}{preàmbul}.
\end{remark}


\begin{exercise}
Totes les estratègies descrites més amunt es fan servir avui dia en escenaris reals, segons les 
característiques, els costos i els compromisos de cadascuna.

Discuteix quina d'aquestes estratègies seria la més adequada per a cadascun d'aquests escenaris:
\begin{itemize}
\item Un sol fabricant dissenya el maquinari i el programari que comunica dos extrems.
Una sola \conceptRef{data line}{línia de dades} s'ha de \conceptRef{multiplex}{multiplexar} per a diverses ordres de control
i també per a diversos fluxos de dades. La prioritat és un ús eficient de l'energia i del buffer (memòria intermèdia).

\item Diversos dispositius comparteixen un \concept{bus}. De tant en tant envien missatges de longituds diferents,
però no un volum enorme de dades. La prioritat és el cost.

\item Diversos dispositius comparteixen un \concept{bus}. Envien contínuament missatges de longituds diferents,
intentant maximitzar la taxa de transmissió efectiva a través del \conceptRef{channel}{canal}. 
La prioritat és el \concept{throughput} (cabal efectiu).
\end{itemize}
\end{exercise}

\section{Si us plau, emplena el formulari}\label{sec:packets:format}

Independentment de l'estratègia emprada, els paquets de dades normalment consten de dues parts: 
\begin{itemize}
 \item El \concept{payload} de dades: la informació que cal transmetre, i
 \item Les \conceptRef{metadata}{metadades}: la metainformació necessària per lliurar el payload.
\end{itemize}
La longitud pot ser un dels \conceptRef{field (packet)}{camps} de metadades inclosos en el paquet, juntament amb 
altres que ajuden a identificar-ne el tipus i la funció.
% 
Els dos extrems de la comunicació han de conèixer la ubicació del payload i de les metadades, així com els camps exactes de metadades inclosos, la seva longitud i el seu significat.
Totes aquestes decisions constitueixen el \conceptRef{packet format}{format del paquet}
o l'\conceptRef{packet structure}{estructura del paquet}.

Una de les maneres més habituals d'organitzar el payload i les metadades és amb un format \conceptRef{header}{capçalera}-\conceptRef{body}{cos}.
En aquest cas, primer s'inclou tota la metainformació, seguida de les dades. 
Altres formats poden incloure un \concept{trailer} o cua final (també conegut com a \concept{tail}, que s'usa sobretot per a la \conceptRef{error detection}{detecció d'errors}
i la correcció), en combinació amb la capçalera o substituint-la.
Totes les configuracions següents són possibles.

\begin{minipage}{\linewidth}
\begin{center}
\begin{bytefield}{32}
\bitbox{11}{Header} & \bitbox{21}{Body}\\[-0.35cm]
\end{bytefield}
% 
\begin{bytefield}{32}
\bitbox{6}{Header} & \bitbox{19}{Body} & \bitbox{7}{Trailer/Tail}\\[-0.35cm]
\end{bytefield}
% 
\begin{bytefield}{32}
\bitbox{21}{Body} & \bitbox{11}{Trailer/Tail}\\[-0.35cm]
\end{bytefield}
% 
\begin{bytefield}{32}
\bitbox{32}{Header}
\end{bytefield}
\end{center}
\end{minipage}


El format del payload d'un paquet normalment el defineix l'usuari. En canvi, el format de la capçalera
està definit estrictament perquè es pugui produir i \conceptRef{parsing}{analitzar} fàcilment.
Aquest format sovint es presenta amb un diagrama que ajuda a identificar les posicions individuals de cada bit i de cada byte,
com en l'exemple (no real) següent:\\

\begin{center}
\begin{bytefield}{32}
\bitheader{0-31} \\
\begin{rightwordgroup}{Header}
\bitbox{16}{Length} & \bitbox{8}{Importance} & \bitbox{4}{Type} & \bitbox{4}{Mode}\\
\bitbox{24}{Destination Address} & \bitbox{8}{(Reserved)} 
\end{rightwordgroup} \\
\begin{rightwordgroup}{Body}
\wordbox[lrt]{1}{Payload} \\
\skippedwords \\
\wordbox[lrb]{1}{} 
\end{rightwordgroup}
\end{bytefield}
\end{center}

\begin{remark}
Els camps també poden tenir longituds fixes o variables, que no necessàriament s'alineen amb els 
\conceptRef{byte boundary}{límits de byte}. Els camps també poden tenir una longitud igual a $1$~bit,
i en aquest cas normalment se'ls anomena \concept{flag} (indicador) o bit de flag.
\end{remark}


\begin{remark}
En diagrames com l'anterior, sovint es fan servir diverses línies (files) perquè es puguin
imprimir i llegir fàcilment. En aquest cas, el significat és el mateix que si tots els camps s'haguessin
presentat en una mateixa línia (fila).
\end{remark}

\begin{exercise}
El contingut de paquet següent s'ha \conceptRef{sniff}{capturat} d'una \conceptRef{network}{xarxa},
expressat en \concept{hexadecimal}. Suposant que el paquet té el format de l'exemple anterior:
% 
\begin{itemize}
\item Indica el valor de cada camp (longitud, importància, tipus, mode, adreça de destinació i payload.
\item Pots endevinar el significat del payload?
\end{itemize}
% 
\begin{center}
\otherBase{00 0c 00 f3 bb bb bb 00 68 65 6c 70}
\end{center}
\end{exercise}

\begin{exercise}
El codi següent accepta dos arguments: (1, \inlineCode{sys.argv[1]}) la \conceptRef{path}{ruta} d'un fitxer d'entrada, 
i (2, \inlineCode{sys.argv[2]}) la \conceptRef{path}{ruta} d'un fitxer de sortida.
Es llegeixen els primers 255 bytes del contingut del fitxer d'entrada, 
es formategen com un paquet, i els bytes d'aquest paquet
s'escriuen al fitxer de sortida.

\begin{center}
\showCode{snippets/simplepacketoutput.py}
\end{center}

\begin{itemize}
\item Descriu el \conceptRef{packet format}{format de paquet} que es fa servir al codi, i les seves limitacions.

\item Amplia el format de paquet de manera que 
 \begin{itemize}
   \item Els paquets puguin tenir més de 255 bytes
   \item Es pugui codificar l'\conceptRef{address}{adreça} del dispositiu de destinació
  \end{itemize}
  
\item Proporciona un diagrama de bytes del nou format, 
així com una explicació del sistema d'adreçament que fas servir.

\item Modifica el codi perquè generi paquets amb el format que proposes.
\end{itemize}
\end{exercise}

\begin{remark}
Llenguatges com Python poden distingir entre fitxers que contenen text i fitxers que contenen dades binàries.
Al codi anterior, els fitxers s'obren en \conceptRef{binary mode}{mode binari} amb els modes \inlineCode{'rb'} i \inlineCode{'wb'}.

Quan s'obren en mode binari, els bytes del fitxer es representen directament amb un objecte \inlineCode{bytes}, un vector 
d'enters entre $0$ i $255$. En \conceptRef{text mode}{mode text}, el mètode \inlineCode{open} processa (\conceptRef{decode}{descodifica}) aquests bytes una mica més,
i retorna una cadena que pot contenir caràcters especials com accents o glifs no llatins. \readMore{https://docs.python.org/3.12/library/functions.html\#open}
\end{remark}

\section{No paren d'arribar}\label{sec:packets:stream}

Quan es desenvolupen aplicacions que fan servir capacitats de xarxa, sovint és útil 
imaginar la comunicació com un \conceptRef{stream}{flux}, \ie, com una canonada conceptual on posem 
les nostres dades (qualsevol quantitat de dades) per un extrem, i surten per l'altre. Si tenim
aquesta capacitat, aleshores és molt més fàcil enviar fitxers de qualsevol mida i transmetre una quantitat
inacabable d'àudio o de vídeo. Malauradament, tot el que tenim per simular aquests fluxos són paquets.

Sovint cal reenviar els \conceptRef{packet}{paquets} a través de diverses xarxes. Els \conceptRef{router}{routers} (encaminadors) no són perfectes
i de vegades estan inoperatius o mal configurats, de manera que no tots els paquets arriben a la seva destinació. A més,
paquets diferents es poden lliurar per rutes diferents que poden tenir longituds diferents, de manera que els paquets poden arribar
desordenats.

Si volem simular un \conceptRef{stream}{flux} de dades, els paquets han de contenir prou metainformació perquè es puguin detectar les pèrdues i 
restaurar l'ordre correcte. Això introdueix una \conceptRef{overhead}{sobrecàrrega} que no és adequada en tots els escenaris --\eg, latència molt baixa--,
de manera que algunes aplicacions basen la seva comunicació de xarxa en \conceptRef{message}{missatges} en lloc de fluxos.

Quan calen fluxos, es fa servir el concepte d'\concept{offset} (desplaçament) per tornar a muntar diversos missatges en 
un sol flux. Cada paquet conté alguns bytes del flux de dades, i l'offset n'indica la ubicació exacta dins d'aquest flux.

\begin{exercise}

El diagrama següent descriu un flux de 48 bytes produït per un node d'origen, 
així com els paquets en què es va dividir abans de la transmissió:\\[-0.25cm]

\begin{center}
\begin{bytefield}[bitheight=2cm]{48}
\bitheader{0-47} \\
\bitbox{10}{\underline{Packet 0} \\Offset 0 \\Length 10}
\bitbox{14}{\underline{Packet 1} \\Offset 10\\Length 14}
\bitbox{6 }{\underline{Packet 2} \\Offset 24\\Length 6 }
\bitbox{18}{\underline{Packet 3} \\Offset 30\\Length 18}
\end{bytefield}
\end{center}

\begin{itemize}
\item Descriu la capçalera mínima que es pot fer servir per transmetre aquest flux.
\item Suposant que el destinatari rep primer el paquet 2, seguit del paquet 0, i que els altres paquets es perden: quins bytes del flux coneix el 
  node de destinació, i quins bytes són desconeguts?
\item Com pot saber el node d'origen que alguns paquets no s'han lliurat?
\item Què passaria si la destinació rebés un missatge fals, amb offset $7$~bytes i longitud $4$~bytes?
\end{itemize}

\end{exercise}



